\section*{引言}
远程状态估计器通过通信网络接收传感器量测，并据此更新状态预测。随机事件触发可依据量测与预测的偏差减少发送次数；在指数二次触发规则下，静默本身也构成估计信息，并可在适当线性高斯模型中保持条件高斯结构\cite{han2015}。

本文关注一个由协议时序产生的矛盾。传感器先依据可信预测反馈决定是否发送，发送位确定后，攻击者才可修改已发送载荷；接收端因此形成偏差，而这个实际预测又会反馈到传感器并影响下一次可信事件。事件记录可能暴露偏差，但审计者看到的统计差异究竟能限制多少真实损伤，并不能仅由“事件可信”直接回答。

我们聚焦三个问题：在攻击峰值与指定记录的总体总变差（total variation，TV）预算同时受限时，精确预测反馈能否对全部允许因果攻击给出真实均方误差（mean-square error，MSE）上界？两项预算如何共同决定最坏损伤的阶数？当实际预测反馈被量化、审计者只保留指定遥测字段时，精确反馈下的零损伤保证会怎样变化？

现有研究已从随机触发估计的创新分布、攻击调度、条件矩约束、可靠信息源和序贯检测等角度研究隐蔽攻击\cite{bai2017,shang2022,guo2025,song2022,xu2023,wei2026}。相关工作中，Shang等提出的 complete Gaussian crater（CGC）分布、历史相关攻击、指定条件下非线性策略的可行性比较，以及可靠信息下的攻击可观测性均有直接成果。表\ref{tab:related}按模型信息、攻击权限和统计约束对比与本文最接近的结果。不同工作的攻击者信息与隐蔽准则不相同；条件矩约束、检测器通过条件和联合记录TV预算不是可互换的指标。

\begin{table}[ht]
\centering\scriptsize
\bicaption{与相关工作的定理级问题对照}{Table \thetable\quad Theorem-level comparison with related work}
\label{tab:related}
\begin{tabularx}{\linewidth}{>{\raggedright\arraybackslash}p{1.55cm}*{4}{>{\raggedright\arraybackslash}X}}
\toprule
工作 & 触发与可信信息 & 攻击及隐蔽约束 & 主要性能结论 & 与本文问题的区别\\
\midrule
Shang等\cite{shang2022} & 源端本地随机事件触发 & 创新/通信通道修改；CGC分布与指定攻击条件 & Theorem 2给协方差递推；Theorem 3及Corollary 3给隐蔽攻击构造 & 受攻远端预测不是触发中心\\
Guo等\cite{guo2025} & 源端创新与攻击调度 & 给定信息下的策略；攻击分支条件均值/协方差约束 & Theorem 1比较给定历史和调度下的条件损伤可行策略，包括非线性策略 & 本文约束指定总体记录分布\\
Wei等\cite{wei2026} & 可靠/可疑信息与序贯检测 & 信息源访问受限；时间及交叉矩条件 & Theorem 2给可行集最坏误差；Corollary 1给特定秩条件 & 触发反馈路径和记录预算不同\\
Li等\cite{li2026} & 前向量化创新及反馈辅助信息 & 量化前创新篡改；指定检测器约束 & Theorem 3优化指定符号攻击；Theorem 4给协方差收敛条件 & 本文量化实际预测反馈中心\\
本文 & 可信发送事件及接收端实际预测反馈 & 发送决定后逐路径有界载荷修改；指定记录TV约束 & 有限时域全策略损伤阶数；量化遥测同律下的残余损伤阶数 & 结论限于本文的协议及记录\\
\bottomrule
\end{tabularx}
\end{table}

本文的区别不在于首次发现可信信息可以约束攻击。与允许攻击者在发送决策前替换预测、或向传感器广播健康预测的模型\cite{song2022,xu2023,wei2026}不同，我们要求反馈忠实于接收端实际预测，事件位不可伪造，载荷修改发生在真实发送决定之后。因而隐藏估计偏差进入下一可信事件的条件概率。结合目标时刻的协方差标签，可校准不同事件历史的健康静默基准。

其关键统计关系将在第\ref{sec:exact}节建立：给定扩充分析历史，下一静默事件的归一化概率等于隐藏偏差的指数二次函数。状态矩阵可逆时，该函数对每个非零偏差严格小于健康值。协方差标签用于控制事件历史概率变化对真实MSE分解中协方差项的影响。由此，峰值限制控制偏差范围，TV预算控制记录可见的总体统计变化，两者共同给出全策略损伤界；构造性下界证明该阶数不能统一改进。

量化反馈扩展用于说明这一保证的适用边界。前向量化观测篡改、量化通信调度以及量化创新下的事件触发估计已有研究\cite{li2026,li2025,shen2025}；本文改变的是进入触发规则的实际预测中心，并明确区分接收端内部信息与审计者记录。量化中心进入触发规则后，多个内部预测可映射为同一反馈码。攻击者可在码格内改变估计而不改变指定的事件及量化遥测记录，因此零TV时仍可能存在正损伤。该结论只针对标量、无饱和量化协议和明确给定的受限记录，不推出完整接收日志下的不可检测性。

本文的理论贡献是上述两个协议下的有限时域损伤界，不是新的触发器、估计器或检测算法。数值部分用于解释真实损伤分解、双预算上下界与量化见证；有限样本不能替代总体定理。本文未证明特定工业系统满足全部信任条件，也未证明受限遥测比完整载荷认证成本更低。全文安排如下：第1节定义系统、攻击权限和记录；第2节给出精确反馈下的损伤机制与定理；第3节分析量化反馈下的残余损伤；第4节报告数值结果；第5节总结适用范围。
